% Companion note to theory Part 0, Definition def:cycle-eps.
% Build: latexmk -pdf -outdir=build product_kernel_gap.tex   (or `make -C theory`)
\documentclass[10pt]{article}
% Shared preamble for the GeoFDI theory notes.
\usepackage[T1]{fontenc}
\usepackage[utf8]{inputenc}
\usepackage{lmodern}
\usepackage[margin=1.1in]{geometry}
\usepackage{amsmath,amssymb,amsthm}
\usepackage{bm}
\usepackage{booktabs}
\usepackage{array}
\IfFileExists{microtype.sty}{\usepackage{microtype}}{}
\usepackage{xcolor}
\usepackage{tikz}
\usetikzlibrary{positioning,arrows.meta}
% todonotes needs TikZ; fall back to a plain red inline note if unavailable so
% the document still builds on a minimal TeX installation.
\IfFileExists{todonotes.sty}{%
  \usepackage[colorinlistoftodos,textsize=footnotesize]{todonotes}%
}{%
  \newcommand{\todo}[2][]{\textbf{\textcolor{red}{[TODO: #2]}}}%
  \newcommand{\listoftodos}{}%
}
% pending-anchor marker usable inside table cells (todonotes floats cannot live inside floats)
\newcommand{\pending}[1]{\textcolor{red!80!black}{[pending: #1]}}
\usepackage[colorlinks=true,linkcolor=blue!55!black,citecolor=blue!55!black,urlcolor=blue!55!black]{hyperref}

% ---------- theorem-like environments ----------
\newtheorem{theorem}{Theorem}[section]
\newtheorem{lemma}[theorem]{Lemma}
\newtheorem{proposition}[theorem]{Proposition}
\newtheorem{corollary}[theorem]{Corollary}
\newtheorem{conjecture}[theorem]{Conjecture}
\theoremstyle{definition}
\newtheorem{definition}[theorem]{Definition}
\newtheorem{example}[theorem]{Example}
\newtheorem{assumption}{Assumption}
\renewcommand{\theassumption}{A\arabic{assumption}}
\theoremstyle{remark}
\newtheorem{remark}[theorem]{Remark}

% ---------- notation macros ----------
\newcommand{\R}{\mathbb{R}}
\newcommand{\Law}{\operatorname{Law}}
\newcommand{\dtv}{d_{\mathrm{TV}}}
\newcommand{\SE}{\mathrm{SE}}
\newcommand{\SO}{\mathrm{SO}}
\newcommand{\Ogrp}{\mathrm{O}}
\newcommand{\eqdist}{\overset{d}{=}}
\newcommand{\gs}{g_{\mathrm{s}}}               % sagittal reflection
\newcommand{\symgrp}{\Sigma}                   % spatio-temporal gait symmetry group
\newcommand{\eps}[1]{\varepsilon_{\mathrm{#1}}}          % per-transition equivariance defect (A1'-A4')
\newcommand{\epsbar}[1]{\bar\varepsilon_{\mathrm{#1}}}   % cycle-level defect (Definition def:cycle-eps)
\newcommand{\Tc}{T_{\mathrm{c}}}                          % closed-loop transition steps per gait cycle
\newcommand{\repQ}{\rho_{Q}}
\newcommand{\repX}{\rho_{X}}
\newcommand{\repU}{\rho_{U}}
\newcommand{\repY}{\rho_{Y}}
\newcommand{\repR}{\rho_{R}}
\newcommand{\repW}{\rho_{W}}                   % block action on the data space W = X x U x Y
\newcommand{\repenv}{\rho_{\mathrm{env}}}      % action on exogenous environment inputs
\newcommand{\push}[1]{{#1}_{\#}}               % pushforward of a measure

\geometry{margin=0.85in}
\setlength{\parskip}{0.25em}
\begin{document}
\noindent\textbf{Why the product-kernel defect does not imply the
$\varepsilon$-robustness lemma}\\
{\small Companion note to GeoFDI theory Part~0, Definition
\emph{cycle-level defects} --- August 15, 2026 (rp001)}
\vspace{0.4em}\hrule\vspace{0.6em}

\noindent\textbf{Claim.}
Two candidate cycle-level definitions of the controller defect are on the
table. (P)~the \emph{product-kernel defect}
$\epsbar{ctrl}^{\mathrm{prod}} := \sup_{x_{0:\Tc-1}} \dtv\bigl(\bigotimes_t
P_{\mathrm{ctrl}}(\cdot\mid x_{t-1},\theta_t),\ \bigotimes_t
P^{\sigma}_{\mathrm{ctrl}}(\cdot\mid x_{t-1},\theta_t)\bigr)$ --- the
equivariance defect of the $\Tc$-fold product kernel with the state
trajectory held \emph{fixed}; and (H)~the \emph{closed-loop hybrid
displacement} of Part~0, eq.~(cycle-eps) --- the total-variation change of
the whole-cycle law when all $\Tc$ controller applications are conjugated
\emph{inside the loop}. Only (H) bounds the lemma's left-hand side: (P) can be
strictly smaller than the closed-loop displacement, so a lemma stated in (P)
would be false. A two-step example exhibits the gap.

\noindent\textbf{Setting.}
$\Tc=2$, binary inputs $u_1,u_2\in\{0,1\}$, and a dynamics with perfect
memory, $x_1=u_1$ ($x_0$ fixed and irrelevant). Nominal and conjugated
controller kernels, with $0<\varepsilon\le\eta<\tfrac12$:
\[
  A_1=\mathrm{Bern}(\tfrac12),\qquad A_1'=\mathrm{Bern}(\tfrac12+\varepsilon),\qquad
  A_2(\cdot\mid x)=\mathrm{Bern}(p_x),\qquad A_2'(\cdot\mid x)=\mathrm{Bern}(p_x+\varepsilon),
\]
with $p_1=1-\eta$ and $p_0=\eta$.
Every per-transition defect equals $\varepsilon$
($\dtv(A_1,A_1')=\dtv(A_2(\cdot\mid x),A_2'(\cdot\mid x))=\varepsilon$ for both
$x$), so $\eps{ctrl}=\varepsilon$ and the coarse linear bound of Part~0 reads
$\epsbar{ctrl}\le\Tc\,\eps{ctrl}=2\varepsilon$.

\noindent\textbf{What the two definitions compute.}
With likelihood ratios $r_t=A_t'/A_t$ and
$D(u_1,x):=\sum_{u_2}A_2(u_2\mid x)\,\lvert 1-r_1(u_1)\,r_2(u_2\mid x)\rvert$,
uniformity of $A_1$ gives
\[
  \text{(P)}\;\; \epsbar{ctrl}^{\mathrm{prod}}=\max_{x\in\{0,1\}}\tfrac14\bigl[D(1,x)+D(0,x)\bigr],
  \qquad
  \text{(H)}\;\; \Delta_{\mathrm{cl}}=\tfrac14\bigl[D(1,x{=}1)+D(0,x{=}0)\bigr],
\]
because in the loop the branch $u_1$ is evaluated at \emph{its own} state
$x_1=u_1$. Enumerating the four outcomes (verified exactly by brute force):
$D(1,1)=4\varepsilon(1-\eta)+4\varepsilon^2$, $D(0,1)=D(1,0)=2\varepsilon$,
$D(0,0)=4\varepsilon(1-\eta)-4\varepsilon^2$, hence
\[
  \epsbar{ctrl}^{\mathrm{prod}}=\tfrac32\varepsilon-\varepsilon\eta+\varepsilon^2
  \qquad\text{versus}\qquad
  \Delta_{\mathrm{cl}}=2\varepsilon(1-\eta)\;>\;\epsbar{ctrl}^{\mathrm{prod}}
  \quad\text{whenever }\eta+\varepsilon<\tfrac12 .
\]
\begin{center}\small
\begin{tabular}{@{}ccccc@{}}
\toprule
$\varepsilon=\eta$ & per-transition $\eps{ctrl}$ & (P) $\epsbar{ctrl}^{\mathrm{prod}}$ & (H) $\Delta_{\mathrm{cl}}$ & linear bound $\Tc\,\eps{ctrl}$ \\
\midrule
$0.10$ & $0.10$ & $0.150$ & $0.180$ & $0.20$ \\
$0.05$ & $0.05$ & $0.075$ & $0.095$ & $0.10$ \\
$0.01$ & $0.01$ & $0.0150$ & $0.0198$ & $0.02$ \\
\bottomrule
\end{tabular}
\end{center}

\noindent\textbf{Why the closed loop exceeds the product defect.}
The quantity $\Delta_{\mathrm{cl}}$ is a \emph{mixture of branch-wise
maxima}, $\tfrac14\sum_{u_1}\max_x D(u_1,x)$, whereas (P) is the
\emph{maximum of the mixture}, $\max_x\tfrac14\sum_{u_1}D(u_1,x)$. The former dominates, and
strictly so when the worst-case conditioning differs across branches: here the step-2
defect that aligns with the step-1 likelihood ratio sits at $x{=}1$ for
$u_1{=}1$ and at $x{=}0$ for $u_1{=}0$, and the loop delivers exactly that
state to each branch. Holding the trajectory fixed in (P) discards precisely
the input--state correlation that makes defects add up. As $\eta\to0$,
$\Delta_{\mathrm{cl}}\to2\varepsilon=\Tc\,\eps{ctrl}$: the coarse linear
bound is attained, consistent with the reviewer's remark that the block-level
TV linear bound is essentially tight. Part~0 therefore defines
$\epsbar{\bullet}$ as the closed-loop hybrid displacement (H), for which the
lemma is a triangle inequality by construction, and treats (P)-type
quantities only as what the state-matched audits estimate per transition.
\end{document}
